\subsection{String Sequence Operations}

A string is a sequence, but it is also a text object with many specialized methods. These methods do not change the original string. Since \texttt{str} is immutable, operations such as searching, splitting, replacing, and case conversion return results while leaving the original object unchanged.

\begin{verbatim}
text = "These pretzels are making me thirsty"

selected_names = [
    "find",
    "index",
    "startswith",
    "endswith",
    "split",
    "join",
    "replace",
    "lower",
    "upper",
    "casefold",
]

print(f"text: {text!r}")
print(f"type(text): {type(text)}")
print()

for name in selected_names:
    print(f"{name:<12} {name in dir(text)}")
\end{verbatim}

Possible output:

\begin{verbatim}
text: 'These pretzels are making me thirsty'
type(text): <class 'str'>

find         True
index        True
startswith   True
endswith     True
split        True
join         True
replace      True
lower        True
upper        True
casefold     True
\end{verbatim}

The methods in this section are ordinary string methods. They are called on a string object using dot notation.

\subsubsection{Searching and Membership: \texttt{in}, \texttt{.find()}, \texttt{.index()}, \texttt{.startswith()}, and \texttt{.endswith()}}

The simplest search operation is \texttt{in}. It returns a Boolean result.

\begin{verbatim}
text = "Nobody expects the Spanish Inquisition"

print(f"{'Spanish' in text = }")
print(f"{'French' in text = }")
print(f"{'Nobody' in text = }")
print(f"{'nobody' in text = }")
\end{verbatim}

Possible output:

\begin{verbatim}
'Spanish' in text = True
'French' in text = False
'Nobody' in text = True
'nobody' in text = False
\end{verbatim}

String search is case-sensitive. \texttt{"Nobody"} and \texttt{"nobody"} are different strings.

The method \texttt{.find()} searches for a substring and returns the index where it first appears. If the substring is not found, it returns \texttt{-1}.

\begin{verbatim}
text = "Nobody expects the Spanish Inquisition"

pos1 = text.find("Nobody")
pos2 = text.find("expects")
pos3 = text.find("Spanish")
pos4 = text.find("French")

print(f"Nobody:  {pos1}")
print(f"expects: {pos2}")
print(f"Spanish: {pos3}")
print(f"French:  {pos4}")
\end{verbatim}

Possible output:

\begin{verbatim}
Nobody:  0
expects: 7
Spanish: 19
French:  -1
\end{verbatim}

The return value is a position. It can be used for slicing.

\begin{verbatim}
text = "Nobody expects the Spanish Inquisition"

pos = text.find("Spanish")

print(f"pos: {pos}")
print(f"text[pos:]: {text[pos:]!r}")
\end{verbatim}

Possible output:

\begin{verbatim}
pos: 19
text[pos:]: 'Spanish Inquisition'
\end{verbatim}

The method \texttt{.index()} is similar to \texttt{.find()}, but it raises an error when the substring is absent.

\begin{verbatim}
text = "No soup for you!"

print(f"text.index('soup'): {text.index('soup')}")

try:
    print(text.index("coffee"))
except ValueError as err:
    print("index search failed")
    print(f"message: {err}")
\end{verbatim}

Possible output:

\begin{verbatim}
text.index('soup'): 3
index search failed
message: substring not found
\end{verbatim}

The choice is practical:

\begin{itemize}
    \item use \texttt{.find()} when absence is an expected result;
    \item use \texttt{.index()} when absence should be treated as an error.
\end{itemize}

The methods \texttt{.startswith()} and \texttt{.endswith()} test boundaries.

\begin{verbatim}
text = "Hello, Newman!"

print(f"{text.startswith('Hello') = }")
print(f"{text.startswith('Newman') = }")
print(f"{text.endswith('Newman!') = }")
print(f"{text.endswith('Jerry') = }")
\end{verbatim}

Possible output:

\begin{verbatim}
text.startswith('Hello') = True
text.startswith('Newman') = False
text.endswith('Newman!') = True
text.endswith('Jerry') = False
\end{verbatim}

These methods are clearer than manual slicing when the intention is boundary testing.

\begin{verbatim}
filename = "report_final.txt"

print(f"manual suffix test: {filename[-4:] == '.txt'}")
print(f"endswith test:      {filename.endswith('.txt')}")
\end{verbatim}

Possible output:

\begin{verbatim}
manual suffix test: True
endswith test:      True
\end{verbatim}

The second line communicates the intention directly: check whether the string ends with \texttt{".txt"}.

\subsubsection{Splitting and Joining: \texttt{.split()} and \texttt{.join()} as Core Text-Sequence Transformations}

The method \texttt{.split()} breaks a string into a list of smaller strings.

\begin{verbatim}
line = "spam eggs bacon"

parts = line.split()

print(f"line:  {line!r}")
print(f"parts: {parts!r}")
print(f"type(parts): {type(parts)}")
\end{verbatim}

Possible output:

\begin{verbatim}
line:  'spam eggs bacon'
parts: ['spam', 'eggs', 'bacon']
type(parts): <class 'list'>
\end{verbatim}

With no argument, \texttt{.split()} splits on runs of whitespace.

\begin{verbatim}
line = "spam   eggs\tbacon\nbeans"

parts = line.split()

print(f"line:  {line!r}")
print(f"parts: {parts!r}")
\end{verbatim}

Possible output:

\begin{verbatim}
line:  'spam   eggs\tbacon\nbeans'
parts: ['spam', 'eggs', 'bacon', 'beans']
\end{verbatim}

A separator can be supplied explicitly.

\begin{verbatim}
csv_line = "Homer,Marge,Bart,Lisa,Maggie"

parts = csv_line.split(",")

print(f"csv_line: {csv_line!r}")
print(f"parts:    {parts!r}")
\end{verbatim}

Possible output:

\begin{verbatim}
csv_line: 'Homer,Marge,Bart,Lisa,Maggie'
parts:    ['Homer', 'Marge', 'Bart', 'Lisa', 'Maggie']
\end{verbatim}

The separator is removed from the result.

The method \texttt{.join()} performs the reverse kind of transformation. It takes an iterable of strings and joins them with a separator string.

\begin{verbatim}
parts = ["Homer", "Marge", "Bart", "Lisa", "Maggie"]

line1 = ",".join(parts)
line2 = " | ".join(parts)
line3 = " ".join(parts)

print(f"line1: {line1!r}")
print(f"line2: {line2!r}")
print(f"line3: {line3!r}")
\end{verbatim}

Possible output:

\begin{verbatim}
line1: 'Homer,Marge,Bart,Lisa,Maggie'
line2: 'Homer | Marge | Bart | Lisa | Maggie'
line3: 'Homer Marge Bart Lisa Maggie'
\end{verbatim}

The separator is the string before \texttt{.join()}.

\begin{verbatim}
parts = ["Ni", "Ni", "Ni"]

text1 = " ".join(parts)
text2 = "-".join(parts)
text3 = "".join(parts)

print(f"text1: {text1!r}")
print(f"text2: {text2!r}")
print(f"text3: {text3!r}")
\end{verbatim}

Possible output:

\begin{verbatim}
text1: 'Ni Ni Ni'
text2: 'Ni-Ni-Ni'
text3: 'NiNiNi'
\end{verbatim}

A common mistake is to try to join non-string elements directly.

\begin{verbatim}
numbers = [10, 20, 12]

try:
    text = ", ".join(numbers)
    print(text)
except TypeError as err:
    print("join failed")
    print(f"message: {err}")
\end{verbatim}

Possible output:

\begin{verbatim}
join failed
message: sequence item 0: expected str instance, int found
\end{verbatim}

The elements must be converted to strings first.

\begin{verbatim}
numbers = [10, 20, 12]

text_parts = []

for nr in numbers:
    text_parts.append(str(nr))

text = " + ".join(text_parts)

print(f"text_parts: {text_parts!r}")
print(f"text:       {text}")
\end{verbatim}

Possible output:

\begin{verbatim}
text_parts: ['10', '20', '12']
text:       10 + 20 + 12
\end{verbatim}

Splitting and joining are central because they move between:

\begin{itemize}
    \item one text sequence;
    \item a list of text components;
    \item another text sequence built from those components.
\end{itemize}

\begin{verbatim}
line = "jerry;elaine;george;kramer"

parts = line.split(";")
pretty = " / ".join(parts)

print(f"line:   {line!r}")
print(f"parts:  {parts!r}")
print(f"pretty: {pretty!r}")
\end{verbatim}

Possible output:

\begin{verbatim}
line:   'jerry;elaine;george;kramer'
parts:  ['jerry', 'elaine', 'george', 'kramer']
pretty: 'jerry / elaine / george / kramer'
\end{verbatim}

\subsubsection{Replacement and Case Transformation: \texttt{.replace()}, \texttt{.lower()}, \texttt{.upper()}, and \texttt{.casefold()}}

The method \texttt{.replace()} creates a new string where occurrences of one substring are replaced with another substring.

\begin{verbatim}
text = "yada yada yada"

new_text = text.replace("yada", "blah")

print(f"text:     {text!r}")
print(f"new_text: {new_text!r}")

print()
print(f"text id:     {id(text)}")
print(f"new_text id: {id(new_text)}")
\end{verbatim}

Possible output:

\begin{verbatim}
text:     'yada yada yada'
new_text: 'blah blah blah'

text id:     140225874656944
new_text id: 140225874657392
\end{verbatim}

The original string is unchanged.

\begin{verbatim}
text = "D'oh! D'oh! D'oh!"

once = text.replace("D'oh", "Woohoo", 1)
twice = text.replace("D'oh", "Woohoo", 2)
all_items = text.replace("D'oh", "Woohoo")

print(f"once:      {once!r}")
print(f"twice:     {twice!r}")
print(f"all_items: {all_items!r}")
\end{verbatim}

Possible output:

\begin{verbatim}
once:      "Woohoo! D'oh! D'oh!"
twice:     "Woohoo! Woohoo! D'oh!"
all_items: "Woohoo! Woohoo! Woohoo!"
\end{verbatim}

The optional third argument limits the number of replacements.

Case transformation methods also return new strings.

\begin{verbatim}
text = "No Soup For You!"

lower_text = text.lower()
upper_text = text.upper()

print(f"text:       {text!r}")
print(f"lower_text: {lower_text!r}")
print(f"upper_text: {upper_text!r}")
\end{verbatim}

Possible output:

\begin{verbatim}
text:       'No Soup For You!'
lower_text: 'no soup for you!'
upper_text: 'NO SOUP FOR YOU!'
\end{verbatim}

These are useful for normalization before comparison.

\begin{verbatim}
answer = "YES"

print(f"answer == 'yes':         {answer == 'yes'}")
print(f"answer.lower() == 'yes': {answer.lower() == 'yes'}")
\end{verbatim}

Possible output:

\begin{verbatim}
answer == 'yes':         False
answer.lower() == 'yes': True
\end{verbatim}

The method \texttt{.casefold()} is a stronger case-normalization method intended for caseless text comparison.

\begin{verbatim}
word1 = "straße"
word2 = "STRASSE"

print(f"word1.lower():    {word1.lower()!r}")
print(f"word2.lower():    {word2.lower()!r}")
print(f"lower same:       {word1.lower() == word2.lower()}")

print()
print(f"word1.casefold(): {word1.casefold()!r}")
print(f"word2.casefold(): {word2.casefold()!r}")
print(f"casefold same:    {word1.casefold() == word2.casefold()}")
\end{verbatim}

Possible output:

\begin{verbatim}
word1.lower():    'straße'
word2.lower():    'strasse'
lower same:       False

word1.casefold(): 'strasse'
word2.casefold(): 'strasse'
casefold same:    True
\end{verbatim}

For simple English text, \texttt{.lower()} is often enough. For broader Unicode text comparison, \texttt{.casefold()} is usually the better tool.

A final example combines several operations in a normal cleaning pipeline.

\begin{verbatim}
line = "  These Pretzels Are Making Me Thirsty!  "

clean = line.strip()
clean = clean.lower()
clean = clean.replace("pretzels", "bagels")

print(f"line:  {line!r}")
print(f"clean: {clean!r}")
\end{verbatim}

Possible output:

\begin{verbatim}
line:  '  These Pretzels Are Making Me Thirsty!  '
clean: 'these bagels are making me thirsty!'
\end{verbatim}

The method \texttt{.strip()} removes leading and trailing whitespace. It is included here only because it commonly appears next to replacement and case normalization in real text-processing code.

The practical rule is:

\begin{quote}
String methods return new string values. Searching reads the string, splitting converts text into pieces, joining builds text from pieces, replacement creates modified text, and case conversion creates normalized text.
\end{quote}